\section{The First Distinction} \label{sec:firstdistinction}

Begin with the one premise. A distinction is made, this and not that. We do not derive why. The question of why there
is something rather than nothing is the single door this account leaves closed, and we will return to it at the end.

\subsection{Why nothing cannot be the whole story}

The premise looks like the weak point. We grant a distinction without deriving it. So it is worth saying exactly in
what sense nothing could not have been the whole of reality, because there is a precise version of this, and there is a
tempting wrong version we are careful not to claim.

The wrong version is dynamical. One often hears that nothing is unstable, that it fluctuates and decays into something.
We do not claim this, and at the foundation we could not. Instability is a process, a change graded by a before and an
after, and the before and after are themselves the integer arrow we only reach in Section~\ref{sec:arrow}, downstream
of the first distinction. To say nothing decays would borrow the very structure we are trying to produce. So we make no
claim that nothing is dynamically unstable. There is no clock yet for it to be unstable in.

The version we can defend is logical. A complete account of reality is a body of distinctions, the things it can tell
apart. The account with no distinction in it is empty. It holds no proposition, no observer, no reference, and in
particular it does not hold the proposition that it is empty. The very sentence ``there is nothing'' draws a
distinction, nothing as against something, and so it steps outside the state it tries to name. There is therefore no
statable, self-containing, distinction-free total state. The blank cannot describe itself, cannot host the claim of its
own blankness, cannot be the whole story, because being the whole story would already be a mark within it.

\begin{principle}[Nothing is not self-containing]
There is no consistent total description that contains no distinction, since any description, including the assertion
of emptiness itself, is already a distinction. Hence the alternatives are exhaustive. Either reality has no describable
content at all, and then there is nothing to explain and no one to explain it, or it has some content, and any content
whatsoever is at least one distinction.
\end{principle}

This is weaker than ``nothing is impossible'' and stronger than ``we simply assume something.'' We assume nothing
beyond the plain fact that there is anything at all, and the existence of this very question is proof of that fact.
From that bare fact a first distinction follows with no further input, since any content is a distinction. The premise
is not an extra ingredient laid on top of existence. It is the minimal reading of existence. Granting it, the rest is
forced.

A distinction has three parts, not two. There is the side that is marked, the side that is not, and the boundary
between them. The marked and the unmarked are mirror images, and the boundary is the absence out of which they are
carved. So the smallest carrier of a distinction is a three-valued quantity, which we write $\{-1, 0, +1\}$. The $0$
is the boundary, the vacuum, the not-yet-marked. The $+1$ and $-1$ are the two sides.

\begin{principle}[The ternary tone]
The minimal value set carrying a distinction is the ternary $\{-1, 0, +1\}$, the unique smallest set possessing both
a vacuum and a nontrivial charge conjugation.
\end{principle}

This is exact, not poetic. Call an alphabet \emph{admissible} if it contains a zero (a vacuum, an absence) and is
closed under a nontrivial sign flip $x \mapsto -x$ (a charge conjugation that genuinely moves at least one element).
The two requirements are forced by what a distinction must do. It must be able to be absent, so a vacuum. And its two
sides must be mirror images of one another, so a sign-flip symmetry. Then the binary $\{0, 1\}$ fails, it has a vacuum
but no mirror partner for $1$. The pair $\{-1, +1\}$ fails, it has a mirror but no vacuum, no place to be empty. The
ternary $\{-1, 0, +1\}$ passes, and it is the smallest set that does. Any larger admissible alphabet, such as
$\{-2,-1,0,1,2\}$, is admissible but not minimal. So three is forced.

The vacuum is special in a way the two sides are not, and this asymmetry matters later. Under the charge conjugation
$C \colon x \mapsto -x$, the sides are exchanged, $+1 \leftrightarrow -1$, while the zero is fixed. So the
structure-preserving symmetries of the tone, the maps that fix the vacuum and commute with $C$, form only the group
$\mathbb{Z}_2$, the identity and the swap, and the three values partition as $1 + 2$, the lone vacuum and the mirrored
pair. We record this because it will distinguish the ternary's three from a different three that appears much later,
the three generations, whose symmetry is the full $S_3$ and whose partition is $3$. The two threes are not the same
object, a fact we verify in Section~\ref{sec:forcedfree}. The ternary remains the irreducible atom of distinction,
with its essential absence.

So from one distinction we have a tone, a three-valued note with a vacuum and a mirror. The next step gives it a time
to sound in.
